\documentclass[a4paper,11pt]{article}
\pdfoutput=1 % if your are submitting a pdflatex (i.e. if you have
    % images in pdf, png or jpg format)
\usepackage{jcappub} % for details on the use of the package, please
                     % see the JCAP-author-manual
\usepackage{aas_macros}
\usepackage[T5]{fontenc} 
\usepackage[utf8]{inputenc}
\usepackage[dvipsnames]{xcolor}
\usepackage{mwe,bm,verbatim,cprotect,xcolor,xspace}
\usepackage{orcidlink}
\usepackage{subfig,csquotes,float, ulem}
\usepackage{bbold}
\usepackage[export]{adjustbox}
\graphicspath{{figs/}}
%\usepackage[allfiguresdraft]{draftfigure}


%%%%%%%%%%%%%%%
%% Reference %%
%%%%%%%%%%%%%%%

\usepackage[nameinlink,noabbrev]{cleveref}
\crefname{equation}{Eq.}{Eqs.}
\crefname{section}{Sec.}{Secs.}
\crefname{figure}{Fig.}{Figs.}
\crefname{table}{Tab.}{Tabs.}
\crefname{appendix}{App.}{Apps.}
\Crefname{figure}{Figure}{Figures}
\Crefname{equation}{Equation}{Equations}
\Crefname{section}{Section}{Sections}
\Crefname{table}{Table}{Tables}
\Crefname{appendix}{Appendix}{Appendices}

%%%%%%%%%%%%%%%%%%%%%%
%%	  Environment   %%
%%%%%%%%%%%%%%%%%%%%%%

\def\be{\begin{equation}}
\def\ee{\end{equation}}
\newcommand{\vs}{\nonumber\\}
\def\ba#1\ea{\begin{align*}#1\end{align*}}
\def\bg#1\eg{\begin{gather}#1\end{gather}}
\newcommand\numberthis{\addtocounter{equation}{1}\tag{\theequation}}

\newcommand{\twofig}[4]{%
 \begin{figure}%
    \centering
    \subfloat{{\includegraphics[width=.4\textwidth]{#1} }}%
    \qquad
    \subfloat{{\includegraphics[width=.4\textwidth]{#2} }}%
    \caption{#4}%
    \label{#3}%
\end{figure}}

\def\switch#1#2{#2}  %1: old 2: revtex4
\newlength{\diagramwidth}
\switch{
\setlength{\diagramwidth}{3.0em}
}{
\setlength{\diagramwidth}{1cm}
}


%%%%%%%%%%%%%%%
%% Reference %%
%%%%%%%%%%%%%%%

\newcommand{\refeq}[1]{Eq.~(\ref{eqn:#1})}          
\newcommand{\refeqs}[2]{Eqs.~(\ref{eqn:#1})--(\ref{eqn:#2})}          
\newcommand{\refEq}[1]{Equation~(\ref{eqn:#1})}          
\newcommand{\reffig}[1]{Fig.~\ref{fig:#1}}  
\newcommand{\reffigs}[2]{Fig.~\ref{fig:#1}--\ref{fig:#2}}  
\newcommand{\refFig}[1]{Figure~\ref{fig:#1}}          
\newcommand{\reftab}[1]{Tab.~\ref{tab:#1}}          
\newcommand{\refsec}[1]{Sec.~\ref{sec:#1}}
\newcommand{\refapp}[1]{App.~\ref{app:#1}}
\newcommand{\specialcell}[2][c]{%
	\begin{tabular}[#1]{@{}c@{}}#2\end{tabular}}

%%%%%%%%%%%%%%
%% Notation %%
%%%%%%%%%%%%%%

\newcommand{\codefont}[1]{{\texttt{#1}}}

\newcommand{\g}{\gamma}
\newcommand{\s}{\sigma}

\def\freeIC{\textsc{freeIC}}
\def\fixedIC{\textsc{fixedIC}}

\newcommand{\fsky}{f_{\rm sky}}
\def\cH{\mathcal{H}}
\def\FMV{\mathcal{F}_{\textsc{mv}}}

\def\Plin{P_{\rm L}}
\def\dlin{\delta_{\rm L}}
\def\knl{k_\text{NL}}
\def\Pnl{P_{\rm m}}

% VECTOR -> bold face

\renewcommand{\emph}[1]{\textit{#1}}

\newcommand{\vn}{\bm{n}}
\newcommand{\nhat}{\hat{n}}
\newcommand{\vnhat}{\bm{\nhat}}
\newcommand{\vx}{\bm{x}}
\newcommand{\vy}{\bm{y}}
\newcommand{\vu}{\bm{u}}
\newcommand{\vr}{\bm{r}}
\newcommand{\vk}{\bm{k}}
\newcommand{\vq}{\bm{q}}
\newcommand{\vp}{\bm{p}}
\newcommand{\Deg}{^{\circ}}

\newcommand{\<}{\langle}
\renewcommand{\>}{\rangle}

\newcommand{\D}{\Delta}
\newcommand{\dLambda}{\delta^{(1)}_\Lambda}
\newcommand{\dLambdazero}{\delta^{(1)}_{\Lambda_0}}
\newcommand{\ddet}{\delta_{\mathrm{det}}}
\newcommand{\ddetLambda}{\delta_{\mathrm{det},\Lambda}}

\newcommand{\eps}{\epsilon}
\newcommand{\Om}{\Omega_m}
\newcommand{\OK}{\Omega_K}
\newcommand{\OL}{\Omega_\Lambda}
\def\rhocr{\rho_{\rm cr,0}}
\newcommand{\rhob}{\bar\rho}
\newcommand{\hubble}{\mathcal{H}}

\def\Mpch{\,h^{-1}\text{Mpc}}
\def\iMpch{\,h\,\text{Mpc}^{-1}}
\def\Mpchcubed{\, h^{-3} \, \text{Mpc}^{3}}
\def\Msunh{\,h^{-1} M_\odot}
\def\Plin{P_\text{L}}

\def\shat{\hat{s}}
\def\epshat{\hat{\eps}}
\def\shatMV{\hat{s}_{\textsc{mv}}}
\def\L{\Lambda}
\def\kmax{k_{\rm max}}
\def\LL{\mathcal{L}}
\def\M{\mathcal{M}}
\def\N{\mathcal{N}}
\def\O{\mathcal{O}}
\def\P{\mathcal{P}}
\def\R{\mathcal{R}}
\def\G{\mathcal{G}}
\def\J{\mathcal{J}}
\def\NGeul{\verb!NGeul!}

\DeclareMathOperator{\Tr}{Tr}
\DeclareMathOperator{\diag}{diag}
\newcommand{\resp}[2]{R^{#1}_{\,\,\,#2}}
\newcommand{\resptwo}[3]{{R_2}^{#1}_{\,\,\,#2,#3}}
\newcommand{\dd}{\mathrm{d}}

\newcommand{\UnitM}{\mathbb{1}}

% bias parameters
\newcommand{\bdelta}{b_{\delta}}
\newcommand{\bdeltatwo}{b_{\delta^2}}
\newcommand{\blapldelta}{b_{\nabla^2\delta}}
\newcommand{\btrMoneMone}{b_{\tr [M^{(1)}M^{(1)}]}}
\newcommand{\bsigmasigma}{b_{\sigma^2}}


% functional derivative
\def\Del{\mathcal{D}}
% tidal derivativ operator
\def\DT{\mathfrak{D}}

% window function / mask
\def\Wi{\mathcal{W}}

% Dirac distribution in nD
\def\dirac#1{\delta_{\mathrm{D}}^{(#1)}}

% notation for third order operator from conv. time derivative
\def\otd{\text{td}}
\def\Otd{O_{\otd}}

% (un)checked boxes
\def\topen{$\square$}
\protected\def\tpass{\textcolor{ForestGreen}{$\checkmark$}}
\def\tfail{\protect\textcolor{red}{$X$}}

% to denote permutations within equations
% - uses 'etoolbox' package
% \perm{n} -> ``$n$ perm.''
% \perm{}  -> ``perm.''
\newcommand{\perm}[1]{ \expandafter\ifstrempty\expandafter{#1} {\mbox{perm.}} {\mbox{$#1$ perm.}} }

\newcommand{\lcdm}{$\Lambda${\rm CDM}}

\DeclareMathOperator{\tr}{tr}
\DeclareMathOperator{\Cov}{Cov}
\let\Re\relax
\DeclareMathOperator{\Re}{Re}
\let\Im\relax
\DeclareMathOperator{\Im}{Im}
\def\lapl{\nabla^2}
\newcommand{\levicivita}{}% initialize
\def\levicivita#1#{\tensor#1{\epsilon}}

\newcommand{\se}{\ensuremath{\sigma_\mathrm{8}}\xspace}

\definecolor{RoyalBlue}{rgb}{0.25,.41,.88}
\definecolor{ForestGreen}{rgb}{.13,.54,.13}
\definecolor{DeepPurple}{rgb}{.45,.1,.75}

\newcommand{\hl}[1]{\textcolor{ForestGreen}{#1}}
\newcommand{\chg}[1]{\textcolor{RoyalBlue}{#1}}
\newcommand{\mn}[1]{\textcolor{DeepPurple}{[\textbf{MN:}~#1]}}
\newcommand{\mnedit}[1]{\textcolor{ForestGreen}{#1}}
\newcommand{\slackcomments}[1]{{\leavevmode\color{blue}[#1]}}

\def\mnras{MNRAS}
\def\jcap{JCAP}
\def\physrep{Phys.~Rep.}

% Use arabic numbering instead of Roman (to avoid mixing with chain label)
\makeatletter
\renewcommand*{\thesection}{\arabic{section}}
\renewcommand*{\thesubsection}{\thesection.\arabic{subsection}}
\renewcommand*{\p@subsection}{}
\makeatother

%%%%%%%%%%%%%%%%
% RESULTS PATH %
%%%%%%%%%%%%%%%%
\newcommand{\figurespath}{figs}


%%%%%%%%%%%%%%%%%%%%%%%%%%%%%
%% EFT cross-cal macros    %%
%%%%%%%%%%%%%%%%%%%%%%%%%%%%%
\newcommand{\PFS}{\text{PFS}}
\newcommand{\DESI}{\text{DESI}}
\newcommand{\fsig}{f\se}
\newcommand{\Mnu}{M_\nu}
\newcommand{\ctilde}{\tilde{c}}
\newcommand{\Pshot}{P_{\rm shot}}
\newcommand{\Bshot}{B_{\rm shot}}
\newcommand{\Ashot}{A_{\rm shot}}
\newcommand{\rsv}{r_{\sigma_v}}
\newcommand{\nbar}{\bar{n}}
\newcommand{\Vov}{V_{\rm overlap}}
\newcommand{\Vfull}{V_{\rm full}}

\title{Data-driven EFT prior calibration from the PFS\,$\times$\,DESI overlap}
\author[a,b,c]{Nhat-Minh Nguyen (Nguyễn Nhật Minh)\orcidlink{0000-0002-2542-7233}}
\affiliation[a]{Center for Data-Driven Discovery, Kavli IPMU (WPI), UTIAS, The University of Tokyo, Kashiwa, Chiba 277-8583, Japan}
\affiliation[b]{Kavli IPMU (WPI), UTIAS, The University of Tokyo, 5-1-5 Kashiwanoha, Kashiwa, Chiba 277-8583, Japan}
\affiliation[c]{Institute For Interdisciplinary Research in Science and Education, ICISE, Quy Nhon, 55121, Vietnam}
\emailAdd{nhat.minh.nguyen@ipmu.jp}
\keywords{Large-scale structure of the universe, Power spectrum, Galaxy surveys, Cosmological parameters from LSS}

\abstract{%
Full-shape analyses of galaxy clustering in the effective field theory
of large-scale structure (EFTofLSS) have emerged as a powerful probe
of cosmological parameters---but marginalizing over \(\sim\!12\) nuisance
parameters per tracer per redshift bin remains the dominant source of
degradation.
%
Recent work has demonstrated that informative priors on these parameters
can substantially sharpen constraints: HOD-informed priors improve
\(\sigma(\se)\) by \(\sim\!23\%\) [Zhang et al.\ 2025], and field-level
simulation-based priors by \(\sim\!50\%\) [Chudaykin, Ivanov \& Philcox 2026].
%
However, all existing approaches rely on halo-occupation-distribution
model assumptions.
%
\emph{Can the nuisance parameters be calibrated directly from the data,
without HOD modelling?}
%
We show that they can.
%
The \(\sim\!1{,}200\;\text{deg}^2\) overlap of PFS-ELG and DESI-ELG at
\(0.8<z<1.6\) provides a natural calibration laboratory: the cross-power
spectrum \(P^{\PFS\times\DESI}_\ell\) carries zero stochastic
contribution, so jointly analysing it with the two auto-spectra cleanly
separates shot noise from signal.
%
This constrains the stochastic parameters \(\Pshot\) and \(a_0\) to
15--45\% of their conservative prior widths---a 2--7\(\times\)
tightening that requires no simulations, no HOD model, and no
assumptions beyond the EFT framework itself.
%
Exporting these calibrated priors to DESI's full
\(14{,}000\;\text{deg}^2\) footprint improves \(\sigma(f\se)\) by
19--35\% and \(\sigma(\Mnu)\) by 48--53\%, recovering 38--63\% of the
gap to perfect nuisance knowledge.
%
The improvement is robust to the assumed FoG velocity ratio
(\(\rsv\) sensitivity: 34--39\%), precisely because the dominant
mechanism is stochastic separation---not counterterm calibration.
%
PFS's lower virial velocity dispersion further motivates an
asymmetric scale cut \(\kmax^{\PFS}>\kmax^{\DESI}\), extending the
usable \(k\)-range and adding \(\sim\!5\%\) on top.
%
The method is model-independent, complementary to simulation-based
priors, and immediately generalizable to any pair of overlapping
spectroscopic surveys---including Euclid\(\,\times\,\)DESI,
Roman\(\,\times\,\)Spec-S5, and DESI-II\(\,\times\,\)MegaMapper.
}


\begin{document}
\date{\today}

\maketitle
\flushbottom

%%% ============================================================
%%% SKELETON SECTIONS — to be filled from SPEC_WRITEUP.md
%%% ============================================================

%-------------------------------------------------------------------
\section{Introduction}
\label{sec:intro}
%-------------------------------------------------------------------

% 1. The problem: EFT nuisance marginalization is expensive
% Cite Chudaykin+ 2025 (prior reparameterization enables CMB+DESI without SNe),
% Zhang+ 2025 (HOD-informed priors improve sigma8 by 23%),
% Chudaykin+ 2026 (field-level SBPs improve H0 by 40%).
% All existing informative-prior approaches depend on HOD assumptions.

% 2. The concept: cross-calibration from survey overlap.
% When two surveys with different target selections observe the same volume,
% their cross-power spectrum constrains per-tracer nuisance parameters
% model-independently.

% 3. PFS x DESI overlap as a concrete realization.

% 4. Paper roadmap.

% TODO(NMN): Write introduction.


%-------------------------------------------------------------------
\section{Data-driven EFT prior calibration}
\label{sec:method}
%-------------------------------------------------------------------

\subsection{The marginalization cost}
\label{sec:marg_cost}

% Review the EFT parameter set (Chudaykin+ 2025, Table I).
% 12 nuisance parameters per tracer per z-bin.
% Which parameters are analytically marginalized and why their priors matter.
% Cite Zhang+ 2025: HOD-informed priors improve sigma(sigma8) by 23%.

% TODO(NMN): Write marginalization cost subsection.


\subsection{Cross-power constraints on nuisance parameters}
\label{sec:cross_power}

% Write the auto-power P^{AA}, P^{BB} and cross-power P^{AB} in terms of
% EFT parameters. Key features:
% - P^{AB} has zero auto-stochastic term (Pshot, a0, a2 enter only in auto).
%   This is the PRIMARY calibration mechanism.
% - Cross-counterterms averaged: (ctr_A + ctr_B)/2.
% - Three spectra at each (k, ell) jointly constrain per-tracer nuisance
%   parameters more tightly than either auto alone.

% TODO(NMN): Write cross-power subsection with equations.


\subsection{Asymmetric scale cuts}
\label{sec:asymmetric_kmax}

% PFS has lower FoG -> EFT converges to higher k -> kmax_PFS > kmax_DESI.
% Formula: kmax_PFS = kmax_DESI x r_sigma_v^{-1/2}.
% At k > kmax_DESI: only P^{AA} and P^{AB} available (6 observables).
% Fisher sums two k-ranges: 9-obs below kmax_DESI + 6-obs above.
% Cite Rubira & Conteddu (2025) for the single-survey precedent.

% TODO(NMN): Write asymmetric kmax subsection.


\subsection{The prior export mechanism}
\label{sec:prior_export}

% Overlap multi-tracer Fisher -> marginalized sigma for each DESI nuisance
% parameter -> Gaussian prior -> added to DESI full-area Fisher.
% Lever arm: V_full / V_overlap ~ 12.

% TODO(NMN): Write prior export subsection.


\subsection{Comparison with simulation-based priors}
\label{sec:sbp_comparison}

% Frame as complementary:
% - SBPs: tight, but dependent on HOD model correctness.
% - Cross-cal: model-independent, but limited by overlap volume.
% - The two can be combined; cross-cal serves as consistency check.
% - Fair comparison: Zhang+ 2025 (~23% on sigma8).
% - Ceiling: Chudaykin+ 2026 (field-level, ~50% on sigma8).

% TODO(NMN): Write SBP comparison subsection.


%-------------------------------------------------------------------
\section{PFS and DESI ELG samples}
\label{sec:surveys}
%-------------------------------------------------------------------

% Survey specs: areas, z-ranges, nbar(z), b1(z).
% Cite DESI Collaboration (arXiv:2503.14738) and Takada+ (2014).
% n(z) from desimodel forecast files (fine-binned, Dz=0.01).
% Volume-weighted rebinning into 4 matched z-bins [0.8,1.0]-[1.4,1.6].

% TODO(NMN): Write survey section with Table 1.

\begin{table}[t]
\centering
\begin{tabular}{lcccccc}
\hline
\multirow{2}{*}{$z$-bin} & \multicolumn{2}{c}{$\nbar$ $[h^3\,\text{Mpc}^{-3}]$} & \multicolumn{2}{c}{$b_1$} & $\Vov$ & $\Vfull^{\DESI}$ \\
 & PFS & DESI & PFS & DESI & \multicolumn{2}{c}{$[h^{-3}\,\text{Gpc}^3]$} \\
\hline
$[0.8,\,1.0]$  & $8.6\times 10^{-4}$ & $8.1\times 10^{-4}$ & 1.26 & 1.24 & 0.59 & 6.9 \\
$[1.0,\,1.2]$  & $9.3\times 10^{-4}$ & $4.6\times 10^{-4}$ & 1.34 & 1.36 & 0.70 & 8.2 \\
$[1.2,\,1.4]$  & $6.4\times 10^{-4}$ & $2.8\times 10^{-4}$ & 1.42 & 1.50 & 0.79 & 9.2 \\
$[1.4,\,1.6]$  & $3.1\times 10^{-4}$ & $1.2\times 10^{-4}$ & 1.50 & 1.65 & 0.85 & 10.0 \\
\hline
\end{tabular}
\caption{Survey properties per matched redshift bin. Number densities are
volume-weighted averages from fine-binned $n(z)$ files.
$\Vov$ uses $1{,}200\;\text{deg}^2$; $\Vfull^{\DESI}$ uses $14{,}000\;\text{deg}^2$.}
\label{tab:surveys}
\end{table}


%-------------------------------------------------------------------
\section{Fisher forecast setup}
\label{sec:fisher}
%-------------------------------------------------------------------

\subsection{EFT parameterization}
\label{sec:eft_params}

% Follow Chudaykin+ 2025 exactly.
% Table 2: fiducial values and prior widths per tracer.

% TODO(NMN): Write EFT param subsection with Table 2.

\begin{table}[t]
\centering
\begin{tabular}{lcccl}
\hline
Parameter & Fiducial & Prior $\mathcal{N}(\mu,\sigma^2)$ & Units & Notes \\
\hline
$b_1\se$      & $b_1\se(z)$ & flat $[0,3]$     & ---              & \\
$b_2\se^2$    & Lazeyras+   & $\mathcal{N}(0,5^2)$  & ---        & \\
$b_{G_2}\se^2$& $-2/7(b_1{-}1)\se^2$ & $\mathcal{N}(0,5^2)$ & --- & \\
$b_{\Gamma_3}$& $23/42(b_1{-}1)$ & $\mathcal{N}(\mu,1^2)$ & --- & \\
$c_0$         & 0           & $\mathcal{N}(0,30^2)$  & $[\Mpch]^2$ & \\
$c_2$         & 30          & $\mathcal{N}(30,30^2)$  & $[\Mpch]^2$ & \\
$c_4$         & 0           & $\mathcal{N}(0,30^2)$  & $[\Mpch]^2$ & \\
$\ctilde$     & 400         & $\mathcal{N}(400,400^2)$ & $[\Mpch]^4$ & FoG \\
$c_1$         & 0           & $\mathcal{N}(0,5^2)$  & $[\Mpch]^2$ & sub-leading \\
$\Pshot$      & 0           & $\mathcal{N}(0,1^2)$  & ---          & rel.\ to $1/\nbar$ \\
$a_0$         & 0           & $\mathcal{N}(0,1^2)$  & ---          & \\
$a_2$         & 0           & $\mathcal{N}(0,1^2)$  & ---          & \\
\hline
\end{tabular}
\caption{EFT nuisance parameters following Chudaykin, Ivanov \& Philcox (2025).
DESI-ELG fiducials shown; PFS-ELG fiducials scaled as described in the text.}
\label{tab:eft_params}
\end{table}


\subsection{PFS-ELG EFT fiducials}
\label{sec:pfs_fiducials}

% Scaling from DESI: ctilde by sigma_v ratio (default 0.75),
% Pshot by nbar, counterterms by b1.
% State the sensitivity parameter r_sigma_v and range tested.

% TODO(NMN): Write PFS fiducials subsection.


\subsection{Multi-tracer Fisher matrix}
\label{sec:mt_fisher}

% Observable vector: P_ell(k) multipoles (ell=0,2,4) for three spectra
% (AA, BB, AB). 9x9 Gaussian covariance at each k from one-loop P(k,mu)
% via Gauss-Legendre. Asymmetric k-range: 9 obs at k <= kmax_DESI,
% 6 obs (AA+AB) at kmax_DESI < k <= kmax_PFS.
% Derivatives via JAX autodiff through ps_1loop_jax.

% TODO(NMN): Write multi-tracer Fisher subsection.


\subsection{Analysis scenarios and benchmarks}
\label{sec:scenarios}

% Define broad, cross-cal, cross-cal-ext, oracle.
% HOD-prior benchmark (Zhang+ 2025): 23% on sigma8, 4% on Omega_m.
% Field-level benchmark (Chudaykin+ 2026): 50% on sigma8.

% TODO(NMN): Write scenarios subsection.


\subsection{Transferability assumption}
\label{sec:transferability}

% DESI's uniform target selection guarantees ELGs in/outside overlap
% have the same EFT parameters.

% TODO(NMN): Write transferability subsection.


%-------------------------------------------------------------------
\section{Results}
\label{sec:results}
%-------------------------------------------------------------------

\subsection{Nuisance calibration in the overlap}
\label{sec:results_overlap}

% -> Fig 1.
% Pshot and a0 tighten most (to 15-45% of broad).
% c_tilde and c0 barely move (~85-95%).
% The cross-spectrum's zero stochastic contribution is the primary mechanism.
% PFS-only does better than DESI-only on c_tilde (asymmetric kmax).

% TODO(NMN): Write overlap calibration results.

\begin{figure}[t]
\centering
\includegraphics[width=\textwidth]{\figurespath/fig1_overlap_calibration.pdf}
\caption{Overlap calibration. Marginalized $\sigma/\sigma_{\rm broad}$ for
selected DESI-ELG nuisance parameters from three Fisher analyses in the
$1{,}200\;\text{deg}^2$ overlap: DESI-only (blue), PFS-only (orange),
and the full PFS$\times$DESI multi-tracer combination (green).
The dashed line marks the broad-prior floor ($\sigma/\sigma_{\rm broad}=1$).
The stochastic parameters $\Pshot$ and $a_0$ show the strongest
tightening---driven by the zero-stochastic property of the cross-spectrum
$P^{\PFS\times\DESI}_\ell$, which cleanly separates shot noise from signal.
The FoG counterterm $\ctilde$ is better constrained by PFS-only than
DESI-only thanks to the asymmetric scale cut
$\kmax^{\PFS}>\kmax^{\DESI}$. In all cases, the multi-tracer analysis
(green) provides the tightest constraint---as expected from information
monotonicity.}
\label{fig:overlap_calibration}
\end{figure}


\subsection{Calibrated priors vs broad}
\label{sec:results_priors}

% -> Fig 2.
% Pshot: nearly order-of-magnitude improvement.
% b2, bG2: modest tightening.
% Counterterms: barely move.

% TODO(NMN): Write calibrated priors results.

\begin{figure}[t]
\centering
\includegraphics[width=\textwidth]{\figurespath/fig2_calibrated_vs_broad.pdf}
\caption{Calibrated vs.\ broad priors. Cross-calibrated prior widths
$\sigma_{\rm cal}$ (colored dots, one per $z$-bin) compared with the
conservative broad priors $\sigma_{\rm broad}$ (light bars) for all
DESI-ELG nuisance parameters. The stochastic parameter $\Pshot$ tightens
by nearly an order of magnitude---this is the parameter that drives the
cosmological improvement in \cref{fig:full_area}. The nonlinear biases
$b_2\se^2$ and $b_{G_2}\se^2$ show modest tightening from one-loop
sensitivity. The counterterms ($c_0$, $c_2$, $c_4$, $\ctilde$) remain
close to their broad values: the overlap data at $k\leq 0.23\iMpch$
carry little information on these $k^2$ and $k^4$ contributions.}
\label{fig:calibrated_vs_broad}
\end{figure}


\subsection{Full-area DESI constraints}
\label{sec:results_full_area}

% -> Fig 3 (money figure).
% Table 3: sigma per scenario, improvement %, efficiency.
% Cross-cal exceeds HOD-prior benchmark for fsigma8 (19% vs 23% line).
% Cross-cal-ext reaches 35%.
% Mnu improves 48-53%.

% TODO(NMN): Write full-area results with Table 3.

\begin{figure}[t]
\centering
\includegraphics[width=\textwidth]{\figurespath/fig3_full_area_constraints.pdf}
\caption{Full-area DESI-ELG constraints. Marginalized $1\sigma$ uncertainties
on $f\se$, $\Mnu$, and $\Om$ from $14{,}000\;\text{deg}^2$, combining
four $z$-bins, under four analysis scenarios.
The cross-calibration at $\kmax=0.20\iMpch$ (orange) already matches the
HOD-prior benchmark of Zhang et al.\ (2025; dashed, 23\% on $\se$);
extending to $\kmax=0.25\iMpch$ (green) surpasses it.
The gap to the oracle (red) is dominated by counterterm parameters that
the $1{,}200\;\text{deg}^2$ overlap cannot constrain at these $k$-scales---the
stochastic parameters, which the cross-calibration \emph{does} constrain,
are no longer the bottleneck.
Dotted line: field-level SBP ceiling (Chudaykin, Ivanov \& Philcox 2026;
50\% on $\se$), which breaks parameter degeneracies inaccessible to
two-point statistics.}
\label{fig:full_area}
\end{figure}

\begin{table}[t]
\centering
\begin{tabular}{lcccccccc}
\hline
Scenario & $\kmax$ & $\sigma(f\se)$ & $\Delta\%$ & $\sigma(\Mnu)$ [eV] & $\Delta\%$ & $\sigma(\Om)$ & $\Delta\%$ \\
\hline
Broad          & 0.20 & 0.076 & ---   & 1.08 & ---   & 0.050 & ---   \\
Cross-cal      & 0.20 & 0.061 & $-19$ & 0.56 & $-48$ & 0.040 & $-19$ \\
Cross-cal-ext  & 0.25 & 0.049 & $-35$ & 0.51 & $-53$ & 0.039 & $-22$ \\
Oracle         & 0.25 & 0.010 & $-87$ & 0.12 & $-89$ & 0.019 & $-61$ \\
\hline
\end{tabular}
\caption{Headline constraints from the full-area DESI-ELG Fisher forecast,
combined across four $z$-bins. Improvement percentages are relative to the
broad-prior baseline.}
\label{tab:headline}
\end{table}


\subsection{Calibration efficiency}
\label{sec:results_efficiency}

% -> Fig 4.
% Efficiency = (sigma_broad - sigma_xcal-ext) / (sigma_broad - sigma_oracle).
% Mnu: ~60%. fsigma8: 38-47% (rises with z). Omega_m: ~35%.

% TODO(NMN): Write efficiency results.

\begin{figure}[t]
\centering
\includegraphics[width=0.65\textwidth]{\figurespath/fig4_calibration_efficiency.pdf}
\caption{Calibration efficiency per $z$-bin, defined as
$(\sigma_{\rm broad}-\sigma_{\rm xcal\text{-}ext})/(\sigma_{\rm broad}-\sigma_{\rm oracle})$.
An efficiency of unity would mean the cross-calibration fully closes the gap
to perfect nuisance knowledge.
$\Mnu$ achieves the highest efficiency ($\sim\!60\%$): the neutrino mass
constraint is most sensitive to precisely the stochastic parameters that
the cross-calibration tightens.
$f\se$ rises with redshift (38\%\,$\to$\,47\%), reflecting PFS's growing
$\nbar$ advantage at higher $z$.
$\Om$ is flat at $\sim\!35\%$---$\Om$ is primarily shape-constrained and
less sensitive to amplitude-related nuisance parameters.}
\label{fig:efficiency}
\end{figure}


\subsection{Sensitivity to the FoG velocity ratio}
\label{sec:results_sensitivity}

% -> Fig 5.
% r_sigma_v in [0.5, 1.0]: sigma(fsigma8) ranges from 0.046 (39%)
% to 0.050 (34%). Confirms asymmetric kmax adds real FoG leverage.
% Even at r_sigma_v = 1.0, 34% improvement from stochastic separation.

% TODO(NMN): Write sensitivity results.

\begin{figure}[t]
\centering
\includegraphics[width=0.65\textwidth]{\figurespath/fig5_sensitivity_rsigmav.pdf}
\caption{Sensitivity to the FoG velocity ratio. Combined $\sigma(f\se)$
from the cross-cal-ext scenario as a function of
$\rsv=\sigma_{v,\PFS}/\sigma_{v,\DESI}$.
Lower $\rsv$ means weaker PFS FoG, hence higher $\kmax^{\PFS}$, more
usable modes, and tighter constraints.
The improvement ranges from 34\% at $\rsv=1$ (identical FoG; the gain
comes entirely from stochastic separation) to 39\% at $\rsv=0.5$ (maximum
FoG leverage via asymmetric $\kmax$).
The mild slope is itself an important result: the cross-calibration is
robust to the assumed $\rsv$, which is currently uncertain, precisely
because the dominant mechanism does not depend on FoG modelling.}
\label{fig:sensitivity}
\end{figure}


%-------------------------------------------------------------------
\section{Discussion}
\label{sec:discussion}
%-------------------------------------------------------------------

% 1. Stochastic dominance: improvement comes from Pshot/a0, not ctilde.
%    This is a strength: robust to FoG modeling.
% 2. Complementarity with SBPs: model-independent cross-check.
% 3. Survey coordination: 19-35% improvement justifies maximizing overlap.
% 4. Generalization: Euclid x DESI, Roman x Spec-S5, DESI-II x MegaMapper.
% 5. Bispectrum extension: would activate c1, break bias degeneracies.
% 6. Caveats: Gaussian Fisher, scaled PFS fiducials, no bispectrum,
%    no window function, approximate SBP comparison.

% TODO(NMN): Write discussion.


%-------------------------------------------------------------------
\section{Conclusions}
\label{sec:conclusions}
%-------------------------------------------------------------------

% Five takeaways:
% 1. EFT nuisance marginalization is the dominant cost in FS analyses.
% 2. Cross-survey overlap provides data-driven, model-independent
%    calibration -- primarily through the zero-stochastic cross-spectrum.
% 3. For PFS x DESI: sigma(fsigma8) improves by 19-35%, sigma(Mnu) by 48-53%.
% 4. Robust to FoG modeling (r_sigma_v sensitivity: 34-39%).
% 5. Complementary to HOD-informed priors; competitive at the Pk level.

% TODO(NMN): Write conclusions.


%-------------------------------------------------------------------
\acknowledgments
%-------------------------------------------------------------------

% TODO(NMN): Add acknowledgments.


%-------------------------------------------------------------------
% Appendices
%-------------------------------------------------------------------
\appendix

\section{Derivative convergence}
\label{app:derivatives}

% JAX autodiff vs 5-point stencil agreement.
% Median relative error < 10% for b1_sigma8, c2, c_tilde.
% Step-size dependence for c_tilde (large fiducial, needs step ~ 50).

% TODO(NMN): Write derivative convergence appendix.


\section{Fisher code verification}
\label{app:fisher_verification}

% Analytic test: P_0(k) = A k^n + 1/nbar reproduces closed-form Fisher.
% MT >= ST (information monotonicity).
% Identical tracers -> F_MT = 2 x F_ST.
% Scenario ordering: sigma_oracle <= sigma_xcal-ext <= sigma_xcal <= sigma_broad.

% TODO(NMN): Write Fisher verification appendix.


%%%%%%%%%%%%%%%%%%%%%%%%%%%%%%%%%%%%%%%%%%%%%%%%%%%%%%%%%%%%%%%%%%%%%%%%%%%
\clearpage
\bibliographystyle{JHEP}
\bibliography{PFSxDESI-multi-fish}

\end{document}
